% !TEX root = ../thesis.tex
% ==============================================================================
%  APPENDIX: SOFTMAX BOTTLENECK IN FEATURE ATTRIBUTION (appendices/appendix_softmax_bottleneck.tex)
% ==============================================================================
%  This appendix provides mathematical proofs and theoretical analysis of
%  why softmax-based attention masks fail in multi-label sentence attribution
%  and justifies the Bernoulli/Sigmoid architectural formulation.
% ==============================================================================

\section{تنگنای سافت‌مکس و نارسایی آن در انتخاب ویژگی‌های علی}\label{app:appendix_softmax_bottleneck}

یکی از چالش‌های بنیادین در طراحی شبکه‌ی مفسر برای متدولوژی استخراج ویژگی‌های علی، پدیده‌ی «فروپاشی تک‌انتخابی»\enfootnote{single-selection collapse} بود. در این پدیده، مدل مفسر در فاز استخراج ویژگی‌ها، به‌طور سیستماتیک از انتخاب بیش از یک جمله امتناع می‌ورزید؛ حتی در سناریوهایی که مجموعه‌ای از گزاره‌ها به‌طور مشترک و مکمل، استدلال نهایی را تشکیل می‌دادند. 

بررسی‌های تحلیلی نشان داد که این رفتار، یک نارسایی در فرآیند بهینه‌سازی یا خطای آموزش نیست، بلکه پیامد قطعی و ریاضیِ استفاده از تابع سافت‌مکس (و معادل مشتق‌پذیر آن، گامبل-سافت‌مکس) برای تولید ماسک اهمیت است. این پیوست در دو بخش، ابتدا محدودیت‌های نظری سافت‌مکس را بر اساس ادبیات پژوهش تبیین می‌کند و سپس با ارائه‌ی اثبات‌های ریاضی، نشان می‌دهد که چگونه استفاده از این تابع در معماری مبتنی بر مکانیزم توجه متقاطع، منجر به تخریب فضای هندسیِ تعبیه‌ها می‌گردد. در نهایت، مبانی ریاضیِ جایگزینی این مکانیزم با توزیع‌های مستقل برنولی (سیگموئید) تشریح می‌شود.


\subsection{محدودیت‌های نظری سافت‌مکس در بازنمایی توزیع‌های چندوجهی}

عملگر سافت‌مکس از نظر بنیادین برای مدل‌سازی توزیع‌های دسته‌بندی با کلاس‌های مانعه‌الجمع طراحی شده است؛ به این معنا که افزایش احتمال یک متغیر، لزوماً به سرکوب احتمال متغیرهای دیگر می‌انجامد (بازی با جمع صفر). با این حال، در پردازش و تولید زبان طبیعی و استدلال، معنا و علیت به ندرت در یک توکن یا گزاره‌ی منفرد متمرکز است. اشتالبرگ و همکاران~\cite{stahlberg2022jam} نشان دادند که مدل‌سازی زبان ذاتاً ماهیتی چندبرچسبی دارد و تحمیل قیدِ $\sum p_i = 1$ در سافت‌مکس، توانایی مدل را در حفظ چندین مسیر معنایی معتبر سرکوب می‌کند. 

علاوه بر قید مجموع واحد، سافت‌مکس از یک محدودیت ساختاری عمیق‌تر به نام «تنگنای سافت‌مکس» رنج می‌برد. یانگ و همکاران~\cite{yang2018breaking} از نظر ریاضی ثابت کردند که ظرفیت تابع سافت‌مکس برای بازنمایی توزیع‌های احتمالی هم‌پوشان، توسط رتبه‌ی ماتریس در فضای پنهان محدود می‌شود. چانگ و مک‌کالوم~\cite{chang2022softmax} نیز نشان دادند که این محدودیت باعث می‌شود مدل‌های زبانی نتوانند توزیع‌های واژگانی چندوجهی را به‌درستی یاد بگیرند. در بستر مدل مفسر ما، زمانی که شبکه باید چندین جمله‌ی حیاتی را به‌طور هم‌زمان استخراج کند، این تنگنای ذاتی مانع از تخصیص امتیاز اهمیت بالا به تمامی آن‌ها می‌شود.


\subsection{فروپاشی هندسی در معماری توجه متقاطع}

علاوه بر محدودیت‌های نظری فوق، تعامل ماسک تولیدشده توسط سافت‌مکس با مکانیزم توجه در شبکه‌ی انرژی مبتنی بر ترنسفور طراحی شده در این پژوهش، باعث فروپاشی ریاضی بازنمایی‌ها می‌گردد. 

فرض کنید مفسر، ۳ جمله را به عنوان پیش‌نیازهای معتبر شناسایی کرده است. به دلیل قید سافت‌مکس، بالاترین امتیازی که شبکه می‌تواند به طور مساوی به این جملات اختصاص دهد، $c \approx 0.33$ است. اعمال این ماسک نرم بر روی بردارهای تعبیه (که پیش‌تر توسط \lr{Sentence-BERT} تولید شده و دارای نُرم $L_2 = 1.0$ هستند)، طول بردارها را به شدت کاهش می‌دهد ($V'_i = c V_i$). این کاهش مقیاس، در بلوک توجه متقاطع شبکه‌ی انرژی، دو خطای سیستماتیک ایجاد می‌کند:

\paragraph{۱. افزایش کاذب دما و کوری مکانیزم توجه.}
در مکانیزم استاندارد توجه، میزان ارتباط میان پرسمان ($Q$) و کلید ($K$) از طریق ضرب داخلی محاسبه می‌شود:
\begin{equation}
    z_i = \frac{Q \cdot K_i^\top}{\sqrt{d_k}}
\end{equation}
هنگامی که بردار کلید توسط ماسک مفسر مقیاس‌دهی می‌شود ($K'_i = c K_i$)، از آنجا که ضرب داخلی یک عملگر خطی است، لوجیت خروجی نیز دقیقاً با همان ضریب مقیاس‌دهی می‌شود ($z'_i = c z_i$). حال اگر این لوجیت جدید را در تابع سافت‌مکس درونیِ لایه‌ی توجه شبکه‌ی انرژی قرار دهیم، داریم:
\begin{equation}
    \text{\lr{Softmax}}(z'_i) = \frac{\exp(c z_i)}{\sum_j \exp(c z_j)} = \frac{\exp(z_i / (1/c))}{\sum_j \exp(z_j / (1/c))}
\end{equation}
این معادله اثبات می‌کند که ضرب‌کردن بردارهای کلید در یک کسرِ $c \in (0, 1)$، از نظر ریاضی کاملاً هم‌ارز با اعمال دمای $\tau = \frac{1}{c}$ در تابع سافت‌مکس است. به عنوان مثال، اختصاص امتیاز $c = 0.33$ معادل با افزایش دما به $\tau \approx 3.0$ است. این افزایش شدید دما، توزیع توجه را به‌شدت مسطح می‌کند. در نتیجه، شبکه‌ی انرژی توانایی خود را در تفکیک جملات مهم از نویز پس‌زمینه از دست داده و عملاً در برابر سیگنال‌ها نابینا می‌گردد.

\paragraph{۲. تخریب دامنه‌ی خروجی در ترکیب ارزش‌ها.}
گام نهایی در لایه‌ی توجه، محاسبه‌ی خروجی ($O$) به صورت میانگین وزن‌دارِ (ترکیب کوژ) بردارهای ارزش ($V$) است که ضرایب وزن‌دهی آن ($\alpha$) از سافت‌مکس به دست می‌آیند و مجموع آن‌ها دقیقاً ۱ است ($\sum \alpha_i = 1$).
\begin{equation}
    O = \sum_{i=1}^{N} \alpha_i V'_i
\end{equation}
از آنجایی که ماسک مفسر، طول بردارهای ارزش را از پیش به $c$ تقلیل داده است ($V'_i = c V_i$)، خروجی مکانیزم توجه به شکل زیر در می‌آید:
\begin{equation}
    O = \sum_{i=1}^{N} \alpha_i (c V_i) = c \sum_{i=1}^{N} \alpha_i V_i
\end{equation}
با توجه به اینکه $\sum \alpha_i V_i$ یک ترکیب محدب از بردارهای یکه است، طول بردار نهایی هرگز نمی‌تواند از مقدار $c$ فراتر رود ($\|O\| \le c$). این در حالی است که شبکه‌ی انرژی، بر روی بردارهای تعبیه‌ی کامل با دامنه‌ی هندسیِ $\sim 1.0$ آموزش دیده است. در نتیجه، این بردار تقلیل‌یافته (با طول $\sim 0.33$)، مدل انرژی را از خمینه‌ی داده‌های توزیع آموزشی خارج کرده و مدل با تلقی آن به عنوان یک داده‌ی \lr{OOD}، انرژی‌های نامنظم و بالایی به آن اختصاص می‌دهد.

برای فرار از تنزل مقیاس و جریمه‌نشدن توسط شبکه‌ی انرژی، بهینه‌ساز مدل یک راه میان‌بر ریاضی می‌یابد، که آن تخصیص تمام بودجه‌ی احتمال $1.0$ به تنها یک جمله و صفر کردن بقیه است. این کار باعث می‌شود تا بردار آن جمله‌ی خاص، بدون کاهش طول و تغییر دما از گلوگاه مکانیزم توجه عبور کند. این پدیده دلیل اصلیِ «تک‌انتخابی شدن» مفسر است.


\subsection{راه‌حل: گذار به توابع برنولی مستقل (سیگموئید)}

برای حل این چالش ساختاری، استفاده از سافت‌مکس با تابع \textbf{سیگموئید} جایگزین شد. تابع سیگموئید هر لوجیت را به‌صورت یک توزیع برنولی مستقل در بازه‌ی $(0, 1)$ نگاشت می‌کند و قید مجموعِ واحد در آن وجود ندارد. در این معماری جدید، اگر شبکه ۳ جمله‌ی حیاتی را شناسایی کند، می‌تواند به طور مستقل به هر سه امتیاز $\sim 1.0$ اختصاص دهد. از طرفی بردارهای تعبیه بدون کاهش مقیاس هندسی از فاز ماسک عبور کرده و وارد مکانیزم توجه متقاطع می‌شوند. در نهایت دمای مکانیزم توجه تغییر نمی‌کند و دامنه‌ی بردارهای خروجی پایدار باقی می‌ماند. البته، آزادسازی شبکه از قید بودجه‌ی سافت‌مکس می‌تواند منجر به انتخاب بی‌رویه و بیش از حد جملات شود. برای کنترل این رفتار، یک عبارت منظم‌سازی خلوت از نوع نُرم $L_1$ به تاب زیان آموزش افزوده شد تا شبکه تنها جملاتی را استخراج کند که حضور آن‌ها حقیقتاً در کاهش انرژی مدل پایین‌دست مؤثر است.
